% Source: https://artofproblemsolving.com/community/c3h1397833_mathcounts_state_level_problem
% Stable identity: aops-topic-1397833-q01
An integer $n$ can be multiplied by a terminating decimal to give $100$ exactly when $\dfrac{100}{n}$ is a terminating decimal. A fraction in lowest terms is a terminating decimal exactly when its denominator has no prime factors other than $2$ and $5$.

The denominator of $\dfrac{100}{n}$ in lowest terms is $\dfrac{n}{\gcd(n,100)}$. If $n$ has a prime factor $p\ne2,5$, then $p$ does not divide $100$, so $p$ survives in this denominator and the decimal does not terminate. If $n=2^a5^b$, the denominator is also of this form and the decimal terminates.

So the good $n$ are the numbers $2^a5^b\le100$:
\[
1,2,4,8,16,32,64,\quad 5,10,20,40,80,\quad 25,50,100,
\]
which is $15$ numbers. The answer is $100-15=\boxed{85}$.
